\documentclass[10pt,letterpaper]{article}
\usepackage[T1]{fontenc}
\usepackage{charter}
\usepackage[margin=0.6in,top=0.5in,bottom=0.55in]{geometry}
\usepackage{microtype,titlesec,enumitem,xcolor,tabularx,fancyhdr}
\usepackage{needspace}
\usepackage[hidelinks]{hyperref}

\definecolor{accent}{HTML}{1F3A5F}
\definecolor{link}{HTML}{1F4E9C}
\newcommand{\lnk}[2]{\href{#1}{\textcolor{link}{#2}}}

\pagestyle{fancy}\fancyhf{}
\renewcommand{\headrulewidth}{0pt}
\cfoot{\footnotesize\color{gray}Roza Akter \textbar{} Page \thepage{} of 2}

\titleformat{\section}{\large\bfseries\scshape\color{accent}}{}{0em}{}[\vspace{-6pt}\color{accent}\rule{\linewidth}{0.6pt}]
\titlespacing*{\section}{0pt}{6pt}{3pt}
\setlist[itemize]{leftmargin=1.1em,itemsep=1pt,topsep=1pt,parsep=0pt,label=\textbullet}
\setlength{\parindent}{0pt}
\linespread{1.0}

% \entry{Left bold}{Right}{Left italic}{Right italic}
\newcommand{\entry}[4]{\textbf{#1}\hfill #2\\ \textit{#3}\hfill\textit{#4}\par}
\newcommand{\proj}[3]{\needspace{4\baselineskip}\textbf{#1}\hfill #3\\ \textit{\small #2}\par}

\begin{document}

\begin{center}
{\LARGE\bfseries\color{accent} Roza Akter}\\[4pt]
\lnk{mailto:rozakterofficial18@gmail.com}{rozakterofficial18@gmail.com} \enspace|\enspace
\lnk{https://www.linkedin.com/in/rozaakter/}{LinkedIn} \enspace|\enspace
\lnk{https://github.com/rozakter18}{GitHub} \enspace|\enspace
\lnk{https://orcid.org/0009-0007-3725-3633}{ORCID}
\end{center}

\section{Education}
\entry{BRAC University, Dhaka, Bangladesh}{Jan 2023 -- Jan 2027 (expected)}{B.Sc.\ in Computer Science and Engineering}{}
\begin{itemize}
  \item \textbf{Relevant coursework:} Bioinformatics (CSE443), Neural Networks (CSE425), Artificial Intelligence (CSE422), Biological Signal and Image Analysis (CSE490F), Algorithms (CSE221), Data Structures (CSE220), Intro.\ to Biology (BIO101), Intro.\ to Biotechnology and Genetic Engineering (BTE101).
\end{itemize}

\section{Research Experience}
\textbf{A Framework for Learning Transferable Representations Across Neurological Disorders}\\
\textbf{Roza Akter}, Sifatul Karim \textbar{} \textit{Undergraduate thesis \textbar{} Supervisor: Dr.\ Swakkhar Shatabda}\hfill Fall 2025 -- Fall 2026 (ongoing)\par
\begin{itemize}
  \item \textbf{Evaluated cross-disease representation transfer between Alzheimer's disease (AD) and Parkinson's disease (PD)}, examining T1-weighted structural MRI, whole-blood transcriptomics and multimodal fusion architectures.
  \item \textbf{Investigated whether models trained on one neurodegenerative disorder leverage shared biological mechanisms to enhance target classification} by systematizing a 4-condition transfer framework across matched ADNI/PPMI cohorts using frozen encoders, patient-disjoint 5-fold CV $\times$ 3 seeds, and paired bootstrap $\Delta$AUC testing.
  \item \textbf{Diagnosed that cross-disease transfer fails to outperform target standalone models}, clarified that apparent transfer signals stem from non-disease confounds (ADNI scanner era, cell-composition shifts), and identified input feature scaling as the primary driver of early-fusion gains.
\end{itemize}

\vspace{3pt}\textbf{Contrastive ECG--Report Alignment on PTB-XL: Measuring the Gap Between Learned Representations and Prompt-Based Diagnosis}\\
Tasbir Ahmed, \textbf{Roza Akter}\\
\textit{AgIA4BioHealth Workshop, IEEE International Conference on Bioinformatics and Biomedicine (BIBM)}, 2026. \textbf{Accepted, to appear}; to be presented Dec 2026.\par
\begin{itemize}
  \item \textbf{Evaluated contrastive 12-lead ECG--language representations} (1D-CNN + Bio\_ClinicalBERT) on 21,799 records, examining prompt-based zero-shot diagnosis against linear probing on frozen embeddings.
  \item \textbf{Investigated whether low zero-shot diagnostic scores stem from weak encoders or prompt misalignment} by systematizing a patient-disjoint 10-fold CV pipeline using standardized English pseudo-reports from 71 SCP descriptions to eliminate InfoNCE false-negative gradients.
  \item \textbf{Diagnosed a 0.334 macro-AUPRC retrieval gap} (0.417 prompt vs.\ 0.751 probe), clarified that text prompts capture only $\sim$32\% of linearly probeable signal, and identified prompt wording as the primary bottleneck in multimodal ECG classification.
\end{itemize}

\section{Research Interests}
\begin{itemize}
  \item \textbf{Multi-omics and transcriptomic analysis} of neurological disease and cancer
  \item \textbf{Representation learning} for biological sequences and proteins
  \item \textbf{Multimodal biomedical machine learning} with confounder-aware, reproducible evaluation
\end{itemize}

\needspace{6\baselineskip}
\section{Publications and Presentations}
\begin{itemize}
  \item Tasbir Ahmed, \textbf{Roza Akter}. ``Contrastive ECG--Report Alignment on PTB-XL: Measuring the Gap Between Learned Representations and Prompt-Based Diagnosis.'' \textit{AgIA4BioHealth Workshop, IEEE International Conference on Bioinformatics and Biomedicine (BIBM)}, 2026. \textbf{Accepted, to appear}; to be presented Dec 2026.
\end{itemize}

\needspace{8\baselineskip}
\section{Selected Research Projects}
\proj{NSCLC Multi-Omics Integration and Knowledge Mining}{MOFA+, GEO}{}
\begin{itemize}
  \item \textbf{Evaluated multi-omics integration} (RNA, methylation, CNV) with MOFA+ across 818 NSCLC patients, examining latent factor associations across tumor stages and validating factor weights on external GEO cohorts.
  \item \textbf{Investigated whether multi-layer molecular programs capture stage-specific tumor progression better than single-omics data} by systematizing a leakage-free 7-step pipeline combining MOFA+ modeling, fold-wise feature selection, 6-resource knowledge mining and GEO factor-weight transfer.
  \item \textbf{Diagnosed a 60\% cohort loss from RPPA inclusion}, clarified a modest bulk stage-predictive signal (CV AUROC 0.603 vs.\ 0.609 for RNA alone), and identified multi-omics-supported candidate drivers validated across clinical cohorts.
\end{itemize}
\proj{Toward Early Alzheimer's Risk Detection in PMOS (PCOS) Patients: A Multi-Tissue, Multi-Method Gene Expression Investigation}{Individual CSE443 research project \textbar{} DESeq2, limma, fgsea, ComBat \textbar{} \lnk{https://github.com/rozakter18/PCOS-AD-Transcriptomic-Investigation}{Code}}{2025}
\begin{itemize}
  \item \textbf{Evaluated individual-level Alzheimer's disease (AD) risk in PCOS patients} across multi-tissue GEO datasets using cross-platform ML classifiers and nearest-centroid metrics.
  \item \textbf{Investigated proposed neuro-endocrine disease links} (insulin resistance, inflammation) by systematizing a pipeline across 4 GEO datasets (brain, adipose, blood) using ComBat harmonization, leakage-free 5-fold CV with fold-wise PCA, and rank-based fgsea pathway analysis.
  \item \textbf{Diagnosed cross-platform scale mismatches causing classifier saturation}, clarified an absence of individual-level risk signal, and identified 14 immune pathways suppressed in PCOS but elevated in AD brain.
\end{itemize}
\proj{The Immune Cells Awake, the Network Sleeps: Rasmussen Encephalitis RNA-seq}{Solo study \textbar{} R, Salmon, tximport, DESeq2, igraph \textbar{} \lnk{https://github.com/rozakter18/Rasmussen-Encephalitis-RNAseq}{Code}}{2026}
\begin{itemize}
  \item \textbf{Evaluated bulk RNA-seq across Rasmussen encephalitis (RE), epilepsy and control brains}, collapsing 98 technical runs into 48 BioSamples to eliminate pseudo-replication.
  \item \textbf{Investigated whether RE alters gene co-expression networks beyond transcript abundance} by systematizing a matched network design (controlling for gene set, edge count, sample size and permutation null) alongside Salmon quantification and Y-linked sex inference.
  \item \textbf{Diagnosed sex-imbalance and outlier confounds in differential expression}, clarified an absence of network topology alterations ($p = 0.34$), and identified a heterogeneous cytotoxic T-cell/NK immune signature (\textit{CXCL9}, \textit{EOMES}).
\end{itemize}
\proj{Benchmarking Transformer Positional Encodings on CRISPR Periodicity}{Single-author study \textbar{} PyTorch, NCBI \textbar{} \lnk{https://github.com/rozakter18/CRISPR-Positional-Encoding-Benchmark}{Code}}{2026}
\begin{itemize}
  \item \textbf{Evaluated four transformer positional encodings} (Sinusoidal, Learned, RoPE, ALiBi) on nucleotide-level CRISPR classification, examining phase-locked boundary error patterns.
  \item \textbf{Investigated whether relative encodings outperform absolute ones on periodic sequences} by systematizing a controlled dataset split across 18 genomes (1,773 units) via 8-mer similarity clustering and equalized training budgets.
  \item \textbf{Diagnosed an 8--10 nt boundary lag} (37--48\% of errors), showed that uncalibrated pipelines distort relative gains by up to 0.28 MCC, and clarified that absolute Sinusoidal and relative RoPE perform equivalently.
\end{itemize}


\proj{Know Thy (Nearest) Neighbor: Retrieval-Augmented Annotation of a Neglected Parasite's Dark Proteome}{Solo project \textbar{} Python, ESM-2, FAISS, Qwen2.5 \textbar{} \lnk{https://github.com/rozakter18/Protein-RAG-Functional-Annotation}{Code}}{2026}
\begin{itemize}
  \item \textbf{Evaluated a retrieval-augmented generation (RAG) framework for protein function prediction}, coupling ESM-2 embeddings (650M parameters) with FAISS vector retrieval and Qwen2.5-7B grounded synthesis on \textit{Leishmania donovani} sequences.
  \item \textbf{Investigated the 32.7\% uncharacterized ``dark proteome'' of visceral leishmaniasis} by systematizing a 5-stage pipeline across 8,887 proteins with ``All-but-the-Top'' anisotropy correction (removing the top 15 PCs), fact-injected LLM prompts and length-ratio filtering ($\le 3\times$).
  \item \textbf{Diagnosed classifier and anisotropy pathologies}, restoring monotonic score calibration (20\% to 98.6\%), clarified tandem gene expansion vs.\ hubness artifacts, and identified 119 validated, high-confidence functional predictions for dark parasite proteins.
\end{itemize}

\section{Technical Skills}
\begin{tabularx}{\linewidth}{@{}l X@{}}
\textbf{Programming} & Python, R, C, C++, SQL \\
\textbf{Bioinformatics} & Salmon, tximport, DESeq2, limma, fgsea, ComBat, MOFA+, igraph, RNA-seq, differential expression, co-expression networks, pathway analysis, Code: \lnk{https://rosalind.info/users/roza18/}{Rosalind} \\
\textbf{Machine learning} & PyTorch, scikit-learn, Transformers, ESM-2, FAISS, contrastive learning, transfer learning, multimodal learning, retrieval-augmented generation \\
\end{tabularx}

\end{document}